% A_methods.tex -- online appendix: the estimators, the implementation and the tuning behind the Methods section
% (Section 5, rehaul/05_methods.tex). Created 2026-10-26 (integration pass, JFEc; METHODS agent). Every item below
% moved here from 05_methods.tex, lightly edited (the previous Methods text is parked, whole, at the end of that file),
% except the subsection "Tuning: provenance and checks" (app:rh_tuning), which is new: the disclosure of decision D14
% (writeup/INTEGRATION_PLAN_2026-10-26.md), with the audit findings correctness-3, fidelity-3, fidelity-2, design-10 and
% correctness-9 (scratchpad audit_2026-10-25.md).
% Wiring: APPENDIX \input's this file (rehaul/A_methods) in the online appendix, best before or after "Implementation
% details" (app:rh_impl, 10_appendix.tex), whose subsections sec:rh_exp_samples, sec:rh_methods_hyper and
% sec:rh_methods_tuning are cited below.
% House rules: every number from a generated macro. \rhPending{} (defined by \providecommand in the number files of the
% preamble; it takes no argument and prints "pending") marks a number whose macro the NUMBERS agent has yet to emit; the
% macro wanted is named in the comment on the same line, with its value in the run files (to be checked by NUMBERS):
%   \rhtuNLateSeries        late-starting series zeroed in every value average of the selection fold (10)
%   \rhtuTreeSearchSpan     the pre-sample of the trees' Optuna search ("1998--2002")
%   \rhtuTreeHarBase        HAR-ladder base of the design the trees were tuned on (2)
%   \rhtuTreeNCols          columns of that design (1{,}144)
%   \rhtuRfBins             histogram bins of the oneDAL forest's splits (256; RULE_rf_skx.txt, max_bins)
%   \rhtuRfMseDiffPct       fold-3 validation MSE of the oneDAL forest relative to the stock forest at W = 3,000,
%                           percent (0.27; RULE_rf_skx.txt REASON, results/rf_engine_bench_2026-10-23)
%   \rhtuAnnFirstYear, \rhtuAnnLastYear   the years of the annual re-tuning check (2014, 2024)
%   \rhtuTreeNSeeds         further seeds of the transferred tree configurations in the 2013 check (4)
% Run files of the checks: results/fomcfix_2026-10-26 (RULE.txt section 7; scores_step4_linear/ for the neighbouring
% penalties and the annual re-tuning, tune2013/ for the trees' 2013 check), on the FOMC-patched final design (D13).

\section{Estimators, implementation and tuning}
\label{app:rh_methods}
% 2026-10-26 (integration, coordinator: the paper keeps its bullet structure): the three prose passages of this file
% outside a list (this opening, the closing sentences of the selection splits, the opening of the tuning subsection)
% and the timing sentences after the nested list of the screens became bold-lead items; wording unchanged.
\begin{itemize}
\item \textbf{Contents.} This appendix gives the estimators of Section~\ref{sec:rh_methods_models} in full, the
  splits on which hyperparameters are selected, the solvers and screens of the penalized regressions, and how the
  frozen hyperparameters of the main text were obtained, with four checks of them.
\end{itemize}

\subsection{Selection splits}
\label{app:rh_methods_splits}
% 2026-10-26 (integration, JFEc): moved from 05_methods.tex ("Hyperparameter-selection splits"); the displays now use
% the lag on the estimation sample, (t,i) \ominus j (audit design-3: the grid lag (t, i-j) of Section 3 steps through
% every grid interval, overnight included, while the final pre-sample fold counts rows of the 13-bar design: fit
% [8202, 47202), gap 325, validation [47527, 50777), results/window_sweep_2014start_2026-10-22/w3000/run_info.json).
\begin{itemize}
\item \textbf{Pre-sample split.} In line with \citet{gu2020empirical}, a split runs forward in time: an estimation
  subsample, an embargo gap of $n_{\mathrm{gap}}$ observations and a validation subsample of $n_{\mathrm{val}}$
  observations. For a split whose validation subsample ends just before $(t,i)$, with $(t,i) \ominus j$ the $j$th
  observation of the estimation sample before $(t,i)$ (Section~\ref{sec:rh_methods_setup}),
  \begin{gather*}
    \mathcal W^{\mathrm{est}}_{t,i} = \{(t,i) \ominus j : j = n_{\mathrm{val}} + n_{\mathrm{gap}} + 1, \dots,
      n_{\mathrm{val}} + n_{\mathrm{gap}} + W\}, \\
    \mathcal W^{\mathrm{gap}}_{t,i} = \{(t,i) \ominus j : j = n_{\mathrm{val}} + 1, \dots,
      n_{\mathrm{val}} + n_{\mathrm{gap}}\}, \qquad
    \mathcal W^{\mathrm{val}}_{t,i} = \{(t,i) \ominus j : j = 1, \dots, n_{\mathrm{val}}\}.
  \end{gather*}
  The estimation subsample is a full window of $W$ observations, so that each candidate is fitted on as many
  observations as a forecast. The linear models choose their penalties once on such a split, \rhorPenSample,
  % 2026-10-28 (author: prune "pooled"). Was: "The pooled linear models choose".
  with $n_{\mathrm{gap}} = \rhdmPlPreEmbargoBars$ bars (\rhdmPlPreEmbargoSess{} trading days) and
  $n_{\mathrm{val}} = \rhdmPlPreFoldBars$ bars (\rhdmPlPreFoldSess{} trading days); the first scored forecast follows
  it. Each candidate penalty $\lambda$ is fitted on $\mathcal W^{\mathrm{est}}_{t,i}$ and scored by the mean squared
  error of $\tilde y$ over $\mathcal W^{\mathrm{val}}_{t,i}$; the minimizer is held fixed at every forecast
  (Appendix~\ref{sec:rh_methods_tuning}; grids and picks in Appendix~\ref{sec:rh_methods_hyper}). The
  closing-interval linear models choose theirs on pre-sample folds of the same form
  (Appendix~\ref{app:rh_closing}).
\item \textbf{Split of the rolling window.} A model that re-selects its hyperparameters during the evaluation, as
  some closing-interval models do, partitions its own window $\mathcal W_{t,i}$, so that the three subsamples together
  hold $W$ observations:
  \[
    \mathcal W^{\mathrm{est}}_{t,i} = \{(t,i) \ominus j : j = n_{\mathrm{val}} + n_{\mathrm{gap}} + 1, \dots, W\},
  \]
  with $\mathcal W^{\mathrm{gap}}_{t,i}$ and $\mathcal W^{\mathrm{val}}_{t,i}$ as above. No forecast of the main text
  re-selects.
\item \textbf{Use of the subsamples.} The target interval $(t,i)$ never enters any of these sets. During a selection,
  a candidate is estimated on $\mathcal W^{\mathrm{est}}_{t,i}$ and evaluated on $\mathcal W^{\mathrm{val}}_{t,i}$;
  every forecast is estimated on the full window $\mathcal W_{t,i}$.
\end{itemize}
% src/models/reclasso_har.py:255-265 (forward_window_split); specs/causal_tune_linear.py:157-162;
% experiments/pooled_linear_rerun_2026_10_19_13.py PRE_FOLD_SESS, PRE_EMBARGO_SESS (moved with the bullet).

\subsection{Penalized linear regressions}
\label{app:rh_methods_linear}
% 2026-10-26 (integration, JFEc): moved from 05_methods.tex ("Penalized linear regressions" and its sub-bullets),
% lightly edited. Changed: the typed "100 times" (audit design-15) -> the weight \rhfiBackbonePf itself; the typed
% "alpha = 0.5" -> "the midpoint of [0, 1]"; the ridge solver (audit fidelity-9: under STANDARDIZE_PENALTY the ridge
% re-solves the standardized system every bar with X'X and X'y kept by rank-one updates, which "replaces the
% Sherman-Morrison update", results/standardize_penalty_2026-10-22/RULE_std.txt items 3 and addendum 3); the timing of
% the screens (fidelity-9 and its verifier: ridge applies the constant, evidence and availability screens every bar,
% the exact-duplicate screen at its cold solves; the lasso and elastic net apply every screen at a cold solve and add
% failing columns between them, the lasso also constant and duplicate ones; the trees have no screen). The
% closing-interval clause of the old screen bullet (hour and overnight indicators) is closing-interval material, dropped
% (audit design-13 (g)). Was: see the parked Methods text at the end of 05_methods.tex.
\begin{itemize}
\item \textbf{Objectives.} Following the penalized HAR formulation of \citet{ZhangZhangCucuringuQian2024}, the
  estimates at origin $(t,i)$ are
  \begin{align*}
    \hat{\bm\beta}^{\mathrm{ridge}}_{t,i} &= \argmin_{\beta_0, \bm\beta} \sum_{(t',i') \in \mathcal W_{t,i}}
      (\tilde y_{t',i'} - \beta_0 - \bm x_{t',i'}^\top \bm\beta)^2 + \lambda \sum_{j=1}^{p} \omega_j\, \beta_j^2, \\
    \hat{\bm\beta}^{\mathrm{enet}}_{t,i} &= \argmin_{\beta_0, \bm\beta} \sum_{(t',i') \in \mathcal W_{t,i}}
      (\tilde y_{t',i'} - \beta_0 - \bm x_{t',i'}^\top \bm\beta)^2
      + \lambda\alpha \sum_{j=1}^{p} \omega_j |\beta_j|
      + \frac{\lambda(1-\alpha)}{2} \sum_{j=1}^{p} \omega_j\, \beta_j^2 ,
  \end{align*}
  where $\lambda > 0$ sets the overall strength of the penalty, $\alpha \in [0,1]$ is the elastic-net mixing
  parameter and $\omega_j \ge 0$ is predictor $j$'s penalty weight; the intercept $\beta_0$ is not penalized. The lasso
  is $\alpha = 1$, and our elastic net fixes $\alpha$ at the midpoint of $[0,1]$. For ridge, $\lambda$ is the grid
  value reported in Table~\ref{tab:rh_dm_hyper}. For the lasso and the elastic net the table reports the
  \texttt{scikit-learn} penalty parameter, and $\lambda$ above is $2n$ times the tabulated value, with $n$ the number
  of observations in the window.
\item \textbf{Penalty weights.} On the full predictor set the HAR-set columns are penalized more lightly than the
  additional predictors:
  \[
    \omega_j =
    \begin{cases}
      \underline\omega = \rhfiBackbonePf & j \in \text{HAR set}, \\
      1 & j \text{ an additional predictor}.
    \end{cases}
  \]
  A penalized regression on the HAR set alone uses $\omega_j = 1$ for every predictor.
\item \textbf{Standardization within the estimation window.} The penalties apply to standardized coefficients
  \citep{Tibshirani1996, FriedmanHastieTibshirani2010}. Each predictor column is centred and scaled by its mean $m_j$
  and standard deviation $s_j$ over $\mathcal W_{t,i}$; the objectives above are solved for the coefficients
  $\gamma_j$ of the standardized columns $(x_j - m_j)/s_j$, with the weights $\omega_j$, and transformed back,
  $\beta_j = \gamma_j / s_j$, the intercept absorbing the centring. Without this step the penalty on a column's
  contribution to the forecast would scale with $1/s_j^2$ for ridge and $1/s_j$ for the lasso, so that columns with
  little spread, such as long-horizon averages, would be penalized far more heavily for no economic reason. The robust
  standardization of Section~\ref{sec:rh_predictors} does not achieve this: it equalizes interquartile ranges rather
  than variances, and its scale factors are not those of the estimation window. For ridge, $m_j$ and $s_j$ are those
  of the current window at every bar; for the lasso and the elastic net they are recomputed at every cold solve of the
  homotopy below and held fixed in between. Solving in standardized coordinates gives the same estimates as penalizing
  the raw columns with weights $\omega_j s_j$ ($\ell_1$) and $\omega_j s_j^2$ ($\ell_2$), and keeps the systems the
  solvers work on well conditioned.
\item \textbf{Screens.} Only a screened subset $\mathcal J_{t,i}$ of the predictors enters a fit:
  \begin{itemize}
  \item columns that are constant within $\mathcal W_{t,i}$, or exact duplicates of an earlier retained column, are
    left out;
  \item a column of an exogenous series enters only if the series' own one-interval column differs from its modal
    value on at least \rhdmModeMinPrints{} intervals of $\mathcal W_{t,i}$, the evidence requirement of the activity
    indicator (Appendix~\ref{app:rh_dp_transform}); a $k$-interval average repeats each observation $k$ times, so its own
    count would overstate the evidence;
  \item an availability average over $k > 1$ intervals is a proportion, with sampling resolution $\sqrt{p(1-p)/k}$
    at its window mean $p$. It is left out when its standard deviation over $\mathcal W_{t,i}$ does not exceed that
    resolution, or when the standard deviation of its difference from an earlier retained availability average of the
    same horizon does not exceed $\sqrt 2$ times the resolution at the pair's mean: its movement is then no larger than
    independent missing observations would produce. These are mostly the long-horizon availability averages of
    regular-hours series, which record the trading calendar rather than the series.
  \end{itemize}
\item \textbf{Timing of the screens.} The HAR set is exempt from the second and third screens. Ridge applies the first screen's constant columns and the
  other two screens at every bar, and the exact duplicates at each of its cold solves (below). The lasso and the
  elastic net apply every screen at each cold solve; between cold solves, a retained column that fails the evidence or
  the availability screen, or, for the lasso, becomes constant or an exact duplicate, is left out at once and the
  solver restarts cold. The tree ensembles use every column.
\item \textbf{Solvers.} With $\bm X_{t,i}$ stacking the rows $(1, \bm x_{t',i'}^\top)$, $(t',i') \in \mathcal W_{t,i}$,
  on the retained columns, ridge solves at every bar the normal equations of the current window in standardized
  coordinates, with the cross-products $\bm X_{t,i}^\top \bm X_{t,i}$ and $\bm X_{t,i}^\top \tilde{\bm y}$ updated by
  rank one as observations enter and leave $\mathcal W_{t,i}$ and recomputed from the current window at fixed
  intervals, its cold solves. The lasso and the elastic net update the active set and
  the coefficients by exact online homotopy \citep{Garrigues2008}, adding the entering observation and dropping the
  leaving one at each step. The homotopy restarts from a cold solve on the current window at fixed intervals; between
  restarts, after each update we check the Karush--Kuhn--Tucker conditions of the current window's problem, and where
  they fail the solution is replaced by a cold solve on the current window, from which the homotopy continues. The OLS
  benchmark is the minimum-norm least-squares solution on the window (Section~\ref{sec:rh_methods_models}).
\end{itemize}
% experiments/close_exogpen_kernel.c:3-12, 650-676 (lambda scales); experiments/pool_linear_pretune.py (STANDARDIZE_PENALTY,
% _std_ridge_solve, min-norm OLS, eigh pseudo-inverse); src/models/reclasso_har.py (WALK_KKT_RESEED, walk_kkt);
% specs/causal_tune_linear.py:118-123, 270-288 (identifiability mask at the tune boundary; between-tune additions for
% lam2 = 0 only); results/standardize_penalty_2026-10-22/RULE_std.txt (addenda 1-3: evidence and availability rules,
% their cadence).

\subsection{Tree ensembles}
\label{app:rh_methods_trees}
% 2026-10-26 (integration, JFEc): moved from 05_methods.tex ("Tree ensembles" and its sub-bullets), lightly edited;
% the forest engine added (audit correctness-9: results/pooled_rf_skx_2014start_w3000/RULE_rf_skx.txt, ENGINE and
% CONFIG; experiments/pooled_trees_rerun_13.py:248-256).
\begin{itemize}
\item \textbf{Gradient boosting} \citep{Friedman2001}, implemented with LightGBM \citep{Ke2017} and XGBoost
  \citep{Chen2016}, minimizes the squared-error loss $\tfrac12 (\tilde y - F)^2$ over boosting rounds $b = 1, \dots, B$
  with learning rate $\eta$. Each regression tree $f_b$ partitions the predictor space into $J$ leaves with leaf weights
  $\bm w$, where leaf $j$ contains the observation subset $\mathcal R_j$ of size $n_j = |\mathcal R_j|$. When the
  $\ell_1$ leaf penalty $\lambda_1^{\mathrm{tree}} = 0$, the additive update and the optimal leaf weight are
  \begin{align*}
    F_b &= F_{b-1} + \eta f_b, \\
    w_j^\ast &= -\frac{\sum_{(t',i') \in \mathcal R_j}\bigl(F_{b-1}(\bm x_{t',i'}) - \tilde y_{t',i'}\bigr)}
      {n_j + \lambda_2^{\mathrm{tree}}} .
  \end{align*}
  At round $b$, tree $f_b$ is grown to minimize the regularized second-order objective
  \[
    \sum_{(t',i') \in \mathcal W_{t,i}} \Bigl[\bigl(F_{b-1}(\bm x_{t',i'}) - \tilde y_{t',i'}\bigr)
      f_b(\bm x_{t',i'}) + \tfrac12 f_b(\bm x_{t',i'})^2\Bigr]
    + \gamma^{\mathrm{tree}} J + \lambda_1^{\mathrm{tree}} \lVert \bm w \rVert_1
    + \tfrac12 \lambda_2^{\mathrm{tree}} \lVert \bm w \rVert_2^2 ,
  \]
  where $\gamma^{\mathrm{tree}}$ penalizes each leaf and $\lambda_1^{\mathrm{tree}}, \lambda_2^{\mathrm{tree}}$
  shrink the leaf weights. Every leaf must contain at least $n_{\min}$ observations ($n_j \ge n_{\min}$), the minimum
  leaf size in LightGBM and, under the unit Hessian of the squared-error loss, the minimum child weight in XGBoost.
\item \textbf{Random forest} \citep{Breiman2001}: the equally weighted average of $B$ regression trees $f_b$, each
  grown on a bootstrap sample of the estimation window with a random subset of the predictors tried at each split,
  \[
    \hat y_{t,i} = \frac1B \sum_{b=1}^{B} f_b(\bm x_{t,i}).
  \]
  The forest is fitted with Intel's oneDAL engine (Intel Extension for Scikit-learn), whose splits search histograms of
  \rhtuRfBins{} bins of each predictor rather than every threshold, % -> \rhtuRfBins
  with the configuration selected for the stock \texttt{scikit-learn} forest (Appendix~\ref{app:rh_tuning}).
\item \textbf{Implementation.} Each method keeps one configuration (Table~\ref{tab:rh_dm_hyper}), refitted at every bar
  on $\mathcal W_{t,i}$ with a fixed random seed. The trees are invariant to monotone transformations of a column, so
  they are not standardized again within the window, and they use every column, with no screen.
\end{itemize}
% Code vs paper (moved): both libraries' L2 objectives are (1/2)(y - F)^2, so h_{t'} = 1 and XGBoost's min_child_weight
% counts rows.

% 2026-10-28 (author: conventional titles). Was: \subsection{Tuning: provenance and checks}
\subsection{Hyperparameter tuning}
\label{app:rh_tuning}
% 2026-10-26 (integration, JFEc): new -- decision D14 (author, 2026-10-26: tuning only needs to put off-the-shelf models
% in the forecasting regime; disclose, do not re-tune). Facts: the selection fold, w3000/frozen_penalties.json (fold 3,
% validation 2013-01-02 .. 2013-12-31) and the late-start rule (src/backtest/executor.py late_start_warmup zeroes rows
% [0, L + W)); results/window_sweep_2014start_2026-10-22/latestart/latestart_spans.txt, fold 3 at w3000: every
% late-start series never live (audit correctness-3, fidelity-3). The tree configurations, cluster/pooltree_lgbm_config.json,
% cluster/xgbtop_xgb_config.json, cluster/pooltree_rf400_config.json (provenance window 24000; experiments/
% tree_menu_dev.py and tree_menu_dev_rf.py on the frozen base-2 panel, src/unification._load_panel, HAR_BASE = 2,
% prep_cache_all_features_b2.npz "300,317 x 1144", masked scaler; audit fidelity-2, design-10). The forest engine,
% RULE_rf_skx.txt and results/rf_engine_bench_2026-10-23 (audit correctness-9). The four checks, results/fomcfix_2026-10-26
% (RULE.txt section 7 (a)-(d), STATUS.md), on the FOMC-patched design. Direction words checked 2026-10-26 against
% scores_step4_linear/tables.md (every-bar full-vs-HAR DM p < 0.001 at all six neighbours and in all three re-tuned
% classes) and tune2013/tune2013_tables.md (forest: T within noise of the best; LightGBM and XGBoost: the best variant,
% doubled learning rate or more rounds, lower on MSE with DM p < 0.01, not on QLIKE, p 0.33 and 0.08).
\begin{itemize}
\item \textbf{Scope of the tuning.} The tuning of the main text only has to put off-the-shelf estimators in the
  forecasting regime; it is not a search for the best configuration, and every comparison of the main text holds the
  hyperparameters fixed. This subsection discloses how the frozen values were obtained and reports four checks of
  them on the final design (Table~\ref{tab:rh_tuning}).
\item \textbf{Linear penalties: the selection fold.} The penalties are selected once, on the split of
  Appendix~\ref{app:rh_methods_splits}: \rhorPenSample. Under the late-start rule of Section~\ref{sec:rh_predictors}, the value
  averages of a series that starts late are zero until a full window of its data exists. In the selection fold this
  holds for all \rhtuNLateSeries{} late-starting series, % -> \rhtuNLateSeries
  VIX, VIX3M, VVIX, the StockTwits series and the volatility-demand series; the first of their averages to
  go live, the one-interval VIX average, does so in \rhlvVixLiveYear{} (Section~\ref{sec:rh_predictors}). The full-set penalties are therefore selected on a design without the implied-volatility series or any alternative
  data, and then govern those columns once they go live. The same holds for the trees' pre-sample check below.
\item \textbf{Tree configurations: transferred from an earlier design.} The configuration of each tree method is the
  best trial of an Optuna search \citep{Akiba2019} (for the random forest, trial \rhdmRfTrial{} of
  \cmTreeProvRfTrials) on the \rhtuTreeSearchSpan{} pre-sample of an earlier design of this panel: % -> \rhtuTreeSearchSpan
  all \rhdmNIntervals{} intervals of the day, overnight included, a HAR ladder of base \rhtuTreeHarBase{}, % -> \rhtuTreeHarBase
  \rhtuTreeNCols{} columns without the FOMC series, % -> \rhtuTreeNCols
  the data before the corrections of Section~\ref{sec:rh_data}, a masked robust scaler and a window of
  \rhdmPooledWindow{} intervals. The configurations are used unchanged on the final design (\rhorNBars{} regular-hours
  bars, \rhorNColsAf{} columns, $W = \rhorWindowBars$), the minimum leaf size $n_{\min}$ included, which is not
  rescaled to the longer window.
\item \textbf{The forest engine.} The forest's configuration was selected with the stock \texttt{scikit-learn} forest,
  which tries every threshold; the forest of the main text is fitted with Intel's oneDAL engine, which splits on
  histograms (Appendix~\ref{app:rh_methods_trees}) and is several times faster. On the selection fold its validation
  MSE is within \rhtuRfMseDiffPct{} percent of the stock forest's at the same window, % -> \rhtuRfMseDiffPct
  and its validation forecasts are nearly perfectly correlated with the stock forest's (Table~\ref{tab:rh_tuning}).
\item \textbf{Check 1: neighbouring penalties.} Each full-set penalized regression is rerun over the whole evaluation
  with its penalty one grid step below and one above the frozen value, the HAR-set forecasts unchanged. At every
  neighbouring penalty the full set's every-bar QLIKE remains significantly below the HAR set's
  (Table~\ref{tab:rh_tuning}).
\item \textbf{Check 2: annual re-tuning.} For each year from \rhtuAnnFirstYear{} to \rhtuAnnLastYear{}, % -> \rhtuAnnFirstYear, \rhtuAnnLastYear
  each penalty is re-selected by the same procedure on a split whose validation subsample is the
  \rhdmPlPreFoldSess{} trading days before the year, using only data before the year, and is applied to the bars of
  that year. From its first live year on, a late-starting series is live in the validation subsample, so the
  re-selected penalties see the implied-volatility series and the alternative data. In every class the full set's
  every-bar gain over the HAR set stays significant (Table~\ref{tab:rh_tuning}).
\item \textbf{Check 3: a pre-sample neighbourhood of each tree configuration.} On the validation year of the
  selection fold only, with a refit at the first bar of every day on the $W$ observations before it, each tree method
  is fitted at its transferred configuration, at \rhtuTreeNSeeds{} further seeds of it, % -> \rhtuTreeNSeeds
  and at variants that move one hyperparameter at a time: the number of leaves (LightGBM) or the depth, the learning
  rate halved and doubled (boosting), half and twice the trees, the minimum leaf size rescaled to the window (for the
  forest also halved), and for the forest the share of predictors tried at a split. For the forest the transferred
  configuration is within noise of the best variant (Diebold--Mariano test of the best variant against it, on the
  daily validation MSE, $p \ge 0.05$). For LightGBM the best variant, a doubled learning rate, and for XGBoost the best
  variant, twice the boosting rounds, have a significantly lower validation MSE than the transferred configuration,
  but not a significantly lower validation QLIKE. This check uses no data after the pre-sample and changes no
  configuration.
\end{itemize}
% The tuning table (Table tab:rh_tuning) is generated by the NUMBERS agent (writeup/generated/rehaul_tuning_table.tex,
% 2026-10-26); this \input needs that file (the Overleaf flattener stops on an \input of a missing file).
% AUTO-GENERATED by writeup/make_rehaul_tuning.py -- do not edit; re-run the script.
% Table tab:rh_tuning (Appendix app:rh_tuning): the forest engine (results/rf_engine_bench_2026-10-23) and
% the tuning checks of results/fomcfix_2026-10-26 (RULE.txt section 7 (a)-(c)).
\begin{table}[htbp]
\centering
\footnotesize
\setlength{\tabcolsep}{3pt}
\begin{threeparttable}
\caption{Tuning provenance and checks}
\label{tab:rh_tuning}
\begin{tabular*}{\linewidth}{@{\extracolsep{\fill}}llrrrrr@{}}
\toprule
\multicolumn{7}{@{}l}{\textit{Panel A: The forest engine (selection fold, validation \rhorPenYears, $W = \rhorWindowDays$ days, full set)}} \\
\addlinespace[2pt]
 & Seconds/refit & MSE & QLIKE & Corr. & & \\
\cmidrule(lr){2-5}
Stock forest & 185 & 0.05329 & 0.1855 & 1 & & \\
oneDAL forest & 32 & 0.05343 & 0.1858 & 0.9982 & & \\
\addlinespace[6pt]
\multicolumn{7}{@{}l}{\textit{Panel B: Linear penalties, full set against the HAR set of the class (\rhorNIntB{} bars and their 16:00 bar)}} \\
\addlinespace[2pt]
 & Penalty & \multicolumn{2}{c}{Every bar} & \multicolumn{2}{c}{16:00} & \\
\cmidrule(lr){3-4}\cmidrule(lr){5-6}
 & & \% vs HAR & DM ($p$) & \% vs HAR & DM ($p$) & \\
\cmidrule(lr){2-6}
Ridge & \ensuremath{10^{3}} (frozen) & \ensuremath{-}6.0 & \ensuremath{-}6.28 (\ensuremath{<}0.001) & +1.8 & 0.67 (0.500) & \\
 & \ensuremath{10^{2}} & \ensuremath{-}6.5 & \ensuremath{-}4.81 (\ensuremath{<}0.001) & +4.1 & 1.21 (0.226) & \\
 & \ensuremath{10^{4}} & \ensuremath{-}5.7 & \ensuremath{-}8.34 (\ensuremath{<}0.001) & \ensuremath{-}1.1 & \ensuremath{-}0.61 (0.544) & \\
 & annual & \ensuremath{-}8.2 & \ensuremath{-}6.40 (\ensuremath{<}0.001) & +1.7 & 0.68 (0.499) & \\
\addlinespace[2pt]
Lasso & \ensuremath{10^{-4}} (frozen) & \ensuremath{-}9.0 & \ensuremath{-}6.34 (\ensuremath{<}0.001) & +4.2 & 1.35 (0.178) & \\
 & \ensuremath{10^{-5}} & \ensuremath{-}8.3 & \ensuremath{-}3.57 (\ensuremath{<}0.001) & +7.3 & 1.85 (0.065) & \\
 & \ensuremath{10^{-3}} & \ensuremath{-}6.9 & \ensuremath{-}8.74 (\ensuremath{<}0.001) & \ensuremath{-}1.5 & \ensuremath{-}0.74 (0.462) & \\
 & annual & \ensuremath{-}6.4 & \ensuremath{-}5.28 (\ensuremath{<}0.001) & \ensuremath{-}1.9 & \ensuremath{-}0.81 (0.418) & \\
\addlinespace[2pt]
Elastic net & \ensuremath{10^{-3}} (frozen) & \ensuremath{-}6.7 & \ensuremath{-}7.36 (\ensuremath{<}0.001) & +0.6 & 0.24 (0.807) & \\
 & \ensuremath{10^{-4}} & \ensuremath{-}8.7 & \ensuremath{-}5.35 (\ensuremath{<}0.001) & +5.3 & 1.54 (0.123) & \\
 & \ensuremath{10^{-2}} & \ensuremath{-}5.6 & \ensuremath{-}10.98 (\ensuremath{<}0.001) & \ensuremath{-}4.8 & \ensuremath{-}3.67 (\ensuremath{<}0.001) & \\
 & annual & \ensuremath{-}6.2 & \ensuremath{-}5.80 (\ensuremath{<}0.001) & \ensuremath{-}2.0 & \ensuremath{-}1.00 (0.319) & \\
\addlinespace[2pt]
\addlinespace[4pt]
\multicolumn{7}{@{}l}{\textit{Panel C: Tree configurations, validation year \rhorPenYears{} only (refit daily)}} \\
\addlinespace[2pt]
 & Best variant & MSE vs $T$ (\%) & $p$ (MSE) & $p$ (QLIKE) & Rank of $T$ & Seeds (\%) \\
\cmidrule(lr){2-7}
LightGBM & learning rate x 2 & \ensuremath{-}2.00 & \ensuremath{<}0.001 & 0.332 & 5 of 8 & \ensuremath{-}0.38 to 0.00 \\
XGBoost & rounds x 2 & \ensuremath{-}2.04 & \ensuremath{<}0.001 & 0.080 & 4 of 8 & 0.00 to +0.32 \\
Random forest & trees x 2 & \ensuremath{-}0.14 & 0.195 & 0.385 & 6 of 9 & \ensuremath{-}0.06 to +0.17 \\
\bottomrule
\end{tabular*}
\begin{tablenotes}[flushleft]
\footnotesize
\item \textit{Notes:} Panel A: the stock \texttt{scikit-learn} random forest, with which the forest's configuration was selected, and Intel's oneDAL forest of the main text (histogram splits, \rhtuRfBins{} bins), fitted with the same configuration on the selection fold of Appendix~\ref{app:rh_methods_splits} (fit on the $W$ bars before a gap, scored on the bars of \rhorPenYears): seconds a refit on 8 threads (AMD EPYC), validation MSE on the adjusted scale and QLIKE, and the correlation of the two forests' validation forecasts. Panel B: each full-set penalized regression at its frozen penalty, at the neighbouring penalties one grid step below and above, and with the penalty re-selected each year from \rhtuAnnFirstYear{} to \rhtuAnnLastYear{} on the \rhdmPlPreFoldSess{} sessions before the year; % QLIKE change against the HAR set of the same class (negative: the full set has the lower loss) and the Diebold--Mariano statistic of the daily average QLIKE (HLN, Newey--West lag \rhorHacLag) with its $p$-value, on the bars of Table~\ref{tab:rh_oos}. Panel C: each tree method at its transferred configuration $T$ and at one-at-a-time variants, refitted at the first bar of every day of \rhorPenYears{} on the $W$ bars before it; the variant with the lowest validation MSE, its MSE relative to $T$, the Diebold--Mariano $p$-values of the daily validation MSE and QLIKE of the best variant against $T$, $T$'s rank among the variants, and the range of $T$'s MSE over \rhtuTreeNSeeds{} further seeds relative to seed 42. Panels B and C use the design with the FOMC corrections of Section~\ref{sec:rh_data}; Panel A the design before them (the corrections change no bar before 2020).
\end{tablenotes}
\end{threeparttable}
\end{table}

